\documentclass[12pt, a4paper]{article}

% PACKAGES
\usepackage[utf8]{inputenc}
\usepackage{amsmath}
\usepackage{amssymb}
\usepackage[margin=1.2in]{geometry}
\usepackage{authblk}
\usepackage{bm}
\newcommand{\lambdabar}{{\mkern0.75mu\mathchar '26\mkern -9.75mu\lambda}}
\linespread{1.2}

\title{A REFLECTION ON HADRONS Part I:  The Phase Dynamics of the Computational Grid}
\author{Haifeng Qiu}

\begin{document}

\begin{center}
    {\uppercase{\Large \textbf{A REFLECTION ON HADRONS}}} \\
    \vspace{0.5em}
    {\large \textit{Part I:  The Phase Dynamics of the Computational Grid}} \\
    \vspace{2em}
    \textbf{Haifeng Qiu} \\
    \today
    \vspace{2em}
\end{center}
\begin{abstract}
This paper further extends the framework of ``Computational Realism'' to the realm of hadron physics, proposing the theory of Grid Phase Dynamics (GPD). This work demonstrates that both strong and weak interactions can be self-consistently unified within the topological fluid evolution of a single non-local nonlinear field equation, without resorting to the abstract postulates of gauge field theories.
\end{abstract}

\section{Complex Nonlinear Non-local Field Equation}

To describe the composite topological soliton systems (hadrons) consisting of fractional-charge sub-standing waves (quarks), we must reshape the dynamics of the underlying vacuum grid by introducing a complex potential field featuring temporal chirality, and assign clear classical physical pictures to a series of topological coupling constants.

\subsection{Complex Bilayer Grid Potential and Temporal Operators}

The underlying computational grid is partitioned into two layers, $b_0$ and $b_1$, where the cross-layer beat delay relationship defines the sign of momentum. We introduce a complex four-potential $A^\mu(\mathbf{x}, t)$ to describe the physical state of this bilayer grid:
\begin{equation}
A^\mu(\mathbf{x}, t) = A^\mu_{b_0}(\mathbf{x}, t) + i A^\mu_{b_1}(\mathbf{x}, t) \label{eq:complex_A}
\end{equation}
Where the imaginary unit $i$ plays a critical role: \textbf{it is not merely an algebraic symbol, but a ``cross-layer information copy latency operator'' representing a $\pi/2$ temporal phase shift between the $b_0$ and $b_1$ layers.}

For wave packets oscillating at high frequencies along the time axis, this inter-layer latency relationship can be quantitatively described via trigonometric combinations:
\begin{enumerate}
    \item \textbf{Electron Sequence ($b_1 \to b_0$ copy latency)}:  
    Information first appears on the $b_1$ layer and is subsequently copied to the $b_0$ layer, meaning $b_0$ lags behind $b_1$. If we set the signal on the $b_1$ layer as $\cos(\omega t)$, then the signal on the $b_0$ layer is $\cos(\omega t - \pi/2) = \sin(\omega t)$. Mathematically, this corresponds to the standard rotating phase:
    \begin{equation}
    A^\mu_{elec} \propto \sin(\omega t) + i \cos(\omega t) = i \left( \cos(\omega t) - i \sin(\omega t) \right) = i e^{-i\omega t} \label{eq:A_elec}
    \end{equation}
    
    \item \textbf{Positron Sequence ($b_0 \to b_1$ copy latency)}:  
    Information first appears on the $b_0$ layer and is subsequently copied to the $b_1$ layer, meaning $b_1$ lags behind $b_0$. If we set the signal on the $b_0$ layer as $\cos(\omega t)$, then the signal on the $b_1$ layer is $\cos(\omega t - \pi/2) = \sin(\omega t)$. This temporal inversion mathematically corresponds to a reversed phase rotation:
    \begin{equation}
    A^\mu_{posi} \propto \cos(\omega t) + i \sin(\omega t) = e^{i\omega t} \label{eq:A_posi}
    \end{equation}
\end{enumerate}

It is straightforward to see that since $(A^\mu_{elec})^* = \sin(\omega t) - i \cos(\omega t) = -i e^{i\omega t} \propto A^\mu_{posi}$, utilizing the complex conjugation transformation $i \to -i$ (which physically corresponds to the inversion of the underlying grid cross-layer temporal sequence), we can strictly map Charge Conjugation Symmetry ($C$-Symmetry) to the temporal sequence inversion symmetry of the underlying grid.

\subsection{Non-local Nonlinear Covariant Field Equation}

The four-dimensional covariant non-local nonlinear field equation governing universal electromagnetic solitons (including both leptons and composite hadrons) is reconstructed as:
\begin{equation}
\exp\left( \frac{c^2}{\omega^2} \square \right) A^\mu = \left[ 1 + \kappa \ln\left(1 + \frac{u_{inv}(\mathbf{x}, t)}{u_{th}(\omega)}\right) \right] A^\mu \label{eq:field_eq}
\end{equation}
The physical meaning and mathematical definition of each term in the equation are as follows:
\begin{enumerate}
    \item \textbf{$\exp(a^2\square)$ Non-local Operator}:  
    Where $a = \bar{\lambda}_c = c/\omega = \frac{\hbar}{mc}$ represents the spatial smoothing scale of the grid (i.e., the reduced Compton wavelength), and $\square = \frac{1}{c^2}\partial_t^2 - \nabla^2$ is the D'Alembertian operator. This exponential operator represents an infinite recursive iteration of information transfer between grid nodes. Its non-local character guarantees the holistic ``particle nature'' of the soliton, naturally introducing a Gaussian-like propagator that fundamentally resolves the point-charge self-energy singularity of classical electrodynamics.
    
    \item \textbf{$\kappa$ Grid Coupling Constant}:  
    This dimensionless constant locks the isotropic characteristic decay depth of the vacuum grid, and its numerical value strictly corresponds to the reciprocal of the Euler's number:
    \begin{equation}
    \kappa \equiv \frac{1}{e_{\text{Euler}}} \approx 0.367879 \label{eq:kappa}
    \end{equation}
    
    \item \textbf{$u_{inv}(\mathbf{x}, t)$ Invariant Lorentz Scalar Energy Density}:  
    In classical electromagnetic field theory, it corresponds to the $T^{00}$ component of the energy-momentum tensor:
    \begin{equation}
    u_{inv} \approx \epsilon_0 \omega^2 A_\nu^* A^\nu \label{eq:u_inv}
    \end{equation}
    In multi-body interactions, this total energy density can be decomposed via algebraic superposition into the soliton's self-energy $u_{self}$, the external background field energy $u_{ext}$, and the coherent constructive interference energy $u_{int}$:
    \begin{equation}
    u_{inv} = u_{self} + u_{ext} + u_{int} \label{eq:u_inv_sum}
    \end{equation}
    where $u_{int} = 2\epsilon_0 \omega^2 \text{Re}(A_{self}^\nu A_{\nu, ext}^*)$.
    
    \item \textbf{$u_{th}(\omega)$ Local Vacuum Energy Threshold}:  
    Represents the critical threshold of the grid medium at which nonlinear shear deformation occurs. Below this threshold, the grid behaves as a linear Maxwell free field ($\square A^\mu = 0$); once the local energy density approaches or exceeds this threshold, the nonlinear self-focusing feedback mechanism is activated:
    \begin{equation}
    u_{th}(\omega) = \Theta \frac{\hbar\omega^4}{c^3} \label{eq:u_th}
    \end{equation}
\end{enumerate}

\subsection{Topological Coupling Constants and Physical Alignment}

To achieve precise alignment between the classical field and macroscopic quantum observables, we must elucidate the geometric origins of the dimensionless constants emerging from Eq. (\ref{eq:field_eq}):
\begin{itemize}
    \item \textbf{$\Theta$ (Absolute Topological Constraint Constant)}:  
    Characterizes the topological invariant of the saturation limit of self-trapping three-dimensional standing waves. It is given by the second-order resonance condition locked by the grid boundary:
    \begin{equation}
    \Theta = \frac{8}{5\pi (e_{\text{Euler}}^2 - 1)} \approx 0.0797138 \label{eq:theta}
    \end{equation}
    This value exhibits a highly precise agreement with the Gaussian integration limit of a three-dimensional ideal continuous manifold, $1/4\pi \approx 0.0795775$.
    
    \item \textbf{$\delta_{grid}$ (Grid Topological Anomaly)}:  
    We define the subtle deviation between the real nonlinear grid boundary $\Theta$ and the ideal three-dimensional manifold boundary $1/4\pi$ as the grid topological anomaly:
    \begin{equation}
    \delta_{grid} \equiv \frac{\Theta - 1/4\pi}{1/4\pi} = 4\pi\Theta - 1 \approx 0.0017129 \label{eq:delta_grid}
    \end{equation}
    Projecting this three-dimensional scalar volume error onto a two-dimensional rotational polarization plane inevitably introduces a polar angle geometric mean factor $\langle\sin^2\theta\rangle = 2/3$. The resulting single-loop torque correction is derived as:
    \begin{equation}
    a_e^{(1)} \approx \frac{2}{3} \delta_{grid} \approx 0.00114 \label{eq:schwinger}
    \end{equation}
    This aligns closely in both physical essence and numerical value with the famous first-order Schwinger radiative correction in Quantum Electrodynamics (QED), $\alpha/2\pi \approx 0.00116$.
    
    \item \textbf{Fluidic Essence of the Fine-Structure Constant $\alpha$}:  
    Substituting the local self-energy of a single soliton at its characteristic cutoff radius $r=a$, $u_{self}(a) = \frac{e^2}{16\pi^2 \epsilon_0 a^4}$, into the logarithmic self-evolution term and dividing by the vacuum deformation threshold $u_{th}$, the characteristic scale $a^4$ undergoes strict algebraic cancellation:
    \begin{equation}
    \frac{u_{self}(a)}{u_{th}(\omega)} = \frac{e^2}{16\pi^2\epsilon_0 a^4} \Big/ \left( \Theta \frac{\hbar c}{a^4} \right) = \frac{1}{4\pi\Theta} \left( \frac{e^2}{4\pi\epsilon_0\hbar c} \right) \equiv \frac{\alpha}{1+\delta_{grid}} \approx \alpha \label{eq:alpha_origin}
    \end{equation}
    This derivation reveals the physical ontology of the fine-structure constant $\alpha$: \textbf{it is by no means an artificially imposed coupling constant, but rather the ``saturation percentage'' of the soliton's self-energy boundary density relative to the absolute vacuum grid breakdown threshold ($u_{th}$)}.
\end{itemize}

With this self-consistent and rigorous non-local nonlinear complex field equation as the physical foundation, we now formally proceed to solve for the topological structures of hadrons.

\section{The Proton Trial Solution and Its Properties}

Under the framework of Grid Phase Dynamics, the proton cannot be viewed simply as a single, structureless symmetric soliton. Instead, it must be described as a composite topological soliton composed of three sub-wavepackets with specific fractionalized charges (the three valence quarks $u_1, u_2, d$), which are spatially separated but highly phase-locked in spacetime, along with an adaptive phase compensation field (the gluon field) required to maintain topological closure.

\subsection{Fractionalized Phase-Frequency Locking and $2\pi$ Topological Closure}

According to the charge-phase mapping rule established in Section 1, the relation between the charge $q$ of a sub-standing wave and its internal relative locking phase difference $\delta$ between forward and backward waves satisfies the linear relationship:
\begin{equation}
\delta_q = \frac{1+q}{2}\pi \label{eq:charge_phase}
\end{equation}
For the three quarks inside the proton, their respective locking phase angles are constrained to:
\begin{itemize}
    \item \textbf{Up Quark ($u_1$)}: $q_{u1} = +2/3 \implies \delta_1 = \frac{1 + 2/3}{2}\pi = \frac{5\pi}{6}$
    \item \textbf{Up Quark ($u_2$)}: $q_{u2} = +2/3 \implies \delta_2 = \frac{1 + 2/3}{2}\pi = \frac{5\pi}{6}$
    \item \textbf{Down Quark ($d$)}: $q_d = -1/3 \implies \delta_3 = \frac{1 - 1/3}{2}\pi = \frac{\pi}{3}$
\end{itemize}

In fluid topology, these three sub-standing waves co-evolve within the same spatial resonant cavity. However, it must be noted that although the sum of the valence quarks' phase-locking angles satisfies the $2\pi$ closure:
\begin{equation}
\sum_{j=1}^3 \delta_j = \delta_1 + \delta_2 + \delta_3 = \frac{5\pi}{6} + \frac{5\pi}{6} + \frac{\pi}{3} = 2\pi \label{eq:sum_phase}
\end{equation}
the linear superposition of their backward wave amplitudes $\sum_{j=1}^3 e^{i\delta_j} = 2e^{i5\pi/6} + e^{i\pi/3} \approx -1.23 + i1.866 \neq -3$. This discrepancy means that pure quark wavepackets cannot achieve perfect out-of-phase cancellation with the forward waves in algebra, failing to form stable, lossless longitudinal standing waves and resulting in severe phase leakage.

To restore local topological closure and prevent phase dissipation, the system must excite an auxiliary ``adaptive phase compensation field.'' \textbf{This is the fluidic origin of the gluon field ($A^\mu_G$)}.

\subsection{Construction of the Proton Composite Fluid Ansatz}

In the fluid-mechanical representation of the underlying grid, we use the complex transverse component $A_{p+} \equiv A_x + i A_y$ to describe the rotational polarization plane of the vector potential. To carry the three internal spatial energy peaks while achieving the macroscopic spin-1/2 fermionic character, we construct the total vector potential $A_{p+}$ as the coherent superposition of the \textbf{quark component $A_{p+, Q}$} and the \textbf{gluon component $A_{p+, G}$} in cylindrical coordinates $(\rho, \theta, z)$:
\begin{equation}
A_{p+}(\rho, \theta, z, t) = \left[ A_{p+, Q}(\rho, \theta, z) + A_{p+, G}(\rho, \theta, z) \right] e^{i\omega_p t} \label{eq:proton_ansatz}
\end{equation}
Where the forward and backward carrier wave envelope functions are defined over the three-dimensional spherical radius $r = \sqrt{\rho^2 + z^2}$ and the grid spatial scale $\bar{\lambda}_p = \frac{\hbar}{m_p c}$:
\begin{equation}
A_{p+, Q} = \sum_{j=1}^3 \left( F_j(r) e^{ik_p z} + F_j(r) e^{i\delta_j} e^{-ik_p z} \right) \Phi_j(\rho, \theta) \label{eq:quark_comp}
\end{equation}
\begin{equation}
A_{p+, G} = F_G(\rho, z) e^{ik_p z} + F_G(\rho, z) e^{i\delta_G(\theta)} e^{-ik_p z} \label{eq:gluon_comp}
\end{equation}
Here, $k_p = \omega_p / c$, and $\Phi_j(\rho, \theta) \approx \phi_0(\rho)\delta(\theta-\theta_j) e^{i\theta_j}$ (where $\theta_j = \{0, 2\pi/3, 4\pi/3\}$) represents the localized angular distribution of the three valence quarks.

To achieve perfect out-of-phase alignment between the total backward wave and the total forward wave (whose fundamental mode amplitude is normalized to $+3$) so that the longitudinal field collapses into a lossless $\sin(k_p z)$ standing wave solution, the projection coefficient $C_G \equiv F_G e^{i\delta_G}$ of the gluon backward wave onto the $n=1$ macroscopic fundamental mode must precisely compensate for the phase distortion generated by the quarks:
\begin{equation}
C_{Q, back} + C_G = -3 \implies C_G = -3 - \sum_{j=1}^3 e^{i\delta_j} \label{eq:gluon_matching}
\end{equation}
Substituting the characteristic phase-locking angles $\delta_j = \{5\pi/6, 5\pi/6, \pi/3\}$ of the proton's quarks, we calculate this complex impedance:
\begin{equation}
C_G = -3 - \left( 2e^{i5\pi/6} + e^{i\pi/3} \right) = -3 - (-1.232 + 1.866i) = -1.768 - 1.866i \approx \mathbf{2.57 e^{-i 2.33 \text{ rad}}} \label{eq:gluon_phase_val}
\end{equation}

\textbf{Physical Alignment Interpretation}: \\
This analytical calculation demonstrates that to eliminate phase leakage, the gluon field is algebraically forced to introduce a \textbf{local phase lag of approximately $-133.46^\circ$ ($-2.33\text{ rad}$)}.

Through this non-linear mode-locking and adaptive phase-pulling, the confinement projection operator $\hat{\mathcal{P}}_{conf}$ reorganizes the spatially separated field sources into a macroscopically lossless, unified proton trial solution:
\begin{equation}
A_{p+, total} \xrightarrow{\hat{\mathcal{P}}_{conf}} -2 F_p(r) \left[ 1 + \varepsilon \cos(3\theta) f(\rho) \right] \sin(k_p z) e^{i(\theta + \omega_p t)} \label{eq:proton_final}
\end{equation}
where $\varepsilon$ is the modulation depth of the three-lobed envelope, and $e^{i(\theta + \omega_p t)}$ is the global rotational phase, ensuring that a $360^\circ$ spatial rotation yields a $4\pi$ topological symmetry with the temporal manifold, imparting the spin-1/2 fermionic property.

\subsection{Transition from Microscopic to Macroscopic: Fourier Harmonic Filtering and Weierstrass Smoothing}

To mathematically prove how the aforementioned microscopic discrete sources converge to a macroscopic smooth modulated wavepacket, we model the angular distribution $\Phi_j(\rho, \theta)$ of the three local wavepackets as $\delta$-functions distributed along a ring of characteristic radius $\rho = \rho_0$:
\begin{equation}
\Phi_j(\rho, \theta) \approx \phi_0(\rho) \delta(\theta - \theta_j) e^{i\beta_j} \label{eq:delta_distribution}
\end{equation}
where $\theta_j = \{0, 2\pi/3, 4\pi/3\}$ are the $C_3$ symmetric angles, and the spatial angular coordinate is chosen as the baseline spin phase $\beta_j = \theta_j$. The microscopic angular distribution $A_\theta^{(micro)}(\theta)$ is expressed as:
\begin{equation}
A_\theta^{(micro)}(\theta) = \sum_{j=0}^2 e^{i\theta_j} \delta(\theta - \theta_j) \label{eq:micro_angular}
\end{equation}
Expanding this into a Fourier series $A_\theta^{(micro)}(\theta) = \sum_{n=-\infty}^\infty c_n e^{in\theta}$, we compute the coefficients $c_n$:
\begin{equation}
c_n = \frac{1}{2\pi} \sum_{j=0}^2 e^{i(1-n)\frac{2\pi j}{3}} \label{eq:fourier_coeff}
\end{equation}
Based on the properties of the sum of roots of unity, $c_n \neq 0$ if and only if $(1-n)$ is an integer multiple of $3$. Thus, the expansion is represented as:
\begin{equation}
A_\theta^{(micro)}(\theta) \propto e^{i\theta} + e^{-2i\theta} + e^{4i\theta} + \dots \label{eq:fourier_series}
\end{equation}

We now apply the infinite-order smoothing operator (heat kernel convolution) $e^{a^2 \nabla^2}$ of the underlying grid. For the angular coordinate $\theta$, the operator reduces to $e^{a^2 \partial_\theta^2}$, which introduces a high-order suppression factor $e^{-a^2 n^2}$ when acting on the Fourier components:
\begin{equation}
A_\theta^{(smoothed)}(\theta) \propto e^{-a^2} e^{i\theta} + e^{-4a^2} e^{-2i\theta} + e^{-16a^2} e^{4i\theta} + \dots \label{eq:smoothed_angular}
\end{equation}
When the grid smoothing scale $a \approx 1$, the fundamental mode $n=1$ has a decay factor of $e^{-1} \approx 0.368$, while the $n=-2$ term is suppressed to $e^{-4} \approx 0.018$ (less than 5\% of the fundamental), and all higher-order terms are algebraically negligible. Retaining only the first two dominant components yields the smoothed distribution:
\begin{equation}
A_\theta^{(smoothed)}(\theta) \approx e^{-a^2} e^{i\theta} \left( 1 + e^{-3a^2} e^{-3i\theta} \right) \label{eq:approx_smoothed}
\end{equation}
Combining the rotating and conjugate terms, this expression simplifies in real space to the smooth envelope form:
\begin{equation}
A_\theta^{(macro)}(\theta) \approx e^{i(\theta + \omega_p t)} \left[ 1 + \varepsilon \cos(3\theta) f(\rho) \right] \label{eq:macro_final}
\end{equation}
where the modulation depth of the three-lobed envelope is $\varepsilon = 2e^{-3a^2} \approx 0.1$. This completes the rigorous algebraic transition from discrete fractional-phase sub-wavepackets to the macroscopic smooth proton trial solution.

\subsection{Algebraic Verification and Order Reduction of the D'Alembertian Operator}

We substitute the trial solution into the D'Alembertian operator $\square = \frac{1}{c^2}\partial_t^2 - \nabla^2$. Applying the Slowly Varying Envelope Approximation (SVEA) from nonlinear wave physics, we neglect the second-order spatial derivatives of the envelope, $\partial_z^2 F_j \approx 0$. The $-k_p^2$ term generated by the second-order time derivative and the $-(-k_p^2)$ term from the second-order spatial derivative \textbf{undergo strict algebraic cancellation}:
\begin{itemize}
    \item \textbf{Action on forward/backward quark components}:
    \begin{equation}
    \square U_{j, Q} \approx \left[ -2ik_p \frac{\partial F_j}{\partial z} - \nabla_\perp^2 F_j \right] e^{ik_p z} \Phi_j e^{i\omega_p t} \label{eq:box_uq}
    \end{equation}
    \begin{equation}
    \square V_{j, Q} \approx \left[ 2ik_p \frac{\partial F_j}{\partial z} - \nabla_\perp^2 F_j \right] e^{i\delta_j} e^{-ik_p z} \Phi_j e^{i\omega_p t} \label{eq:box_vq}
    \end{equation}
    
    \item \textbf{Action on forward/backward gluon components}:
    \begin{equation}
    \square U_G \approx \left[ -2ik_p \frac{\partial F_G}{\partial z} - \nabla_\perp^2 F_G \right] e^{ik_p z} e^{i\omega_p t} \label{eq:box_ug}
    \end{equation}
    \begin{equation}
    \square V_G \approx \left[ 2ik_p \frac{\partial F_G}{\partial z} - \nabla_\perp^2 \left( F_G e^{i\delta_G(\theta)} \right) e^{-i\delta_G(\theta)} \right] e^{i\delta_G(\theta)} e^{-ik_p z} e^{i\omega_p t} \label{eq:box_vg}
    \end{equation}
\end{itemize}

Substituting these into the total nonlinear field equation and extracting the co-frequency components yields a set of first-order, one-dimensional spatial differential equations, namely the \textbf{multibody Coupled-Mode Equations (CMEs)} inside the proton:
\begin{itemize}
    \item \textbf{Forward wave coupling channel}:
    \begin{equation}
    \sum_{j=1}^3 \left( -2ik_p \frac{\partial F_j}{\partial z} - \nabla_\perp^2 F_j \right) \Phi_j + \left( -2ik_p \frac{\partial F_G}{\partial z} - \nabla_\perp^2 F_G \right) = \frac{g(u)}{a^2} \left( \sum_{j=1}^3 F_j \Phi_j + F_G \right) \label{eq:cme_forward}
    \end{equation}
    
    \item \textbf{Backward wave coupling channel}:
    \begin{equation}
    \begin{split}
    \sum_{j=1}^3 \left( 2ik_p \frac{\partial F_j}{\partial z} - \nabla_\perp^2 F_j \right) e^{i\delta_j} \Phi_j + \left( 2ik_p \frac{\partial F_G}{\partial z} - \nabla_\perp^2 \left( F_G e^{i\delta_G(\theta)} \right) e^{-i\delta_G(\theta)} \right) e^{i\delta_G(\theta)} \\
= \frac{g(u)}{a^2} \left( \sum_{j=1}^3 F_j e^{i\delta_j} \Phi_j + F_G e^{i\delta_G(\theta)} \right) \label{eq:cme_backward}
    \end{split}
    \end{equation}
\end{itemize}

\textbf{Strict cancellation of imaginary terms and topological locking}: In the backward channel equation, the transverse Laplacian acting on the gluon angular phase $\delta_G(\theta)$ excites an imaginary current containing the spatial gradient (topological phase curvature):
\begin{equation}
\text{Im} \left[ \nabla_\perp^2 \left( F_G e^{i\delta_G} \right) e^{-i\delta_G} \right] = 2(\nabla_\perp F_G) \cdot (\nabla_\perp \delta_G) + F_G \nabla_\perp^2 \delta_G \label{eq:im_current}
\end{equation}
This term acts as a \textbf{``topological sponge''}, which, by adjusting the spatial gradient of the gluon phase, \textbf{completely absorbs and cancels all non-aligned imaginary residuals generated by the quarks' fractional phases during interference}. Consequently, the net imaginary leakage at the outer boundary of the proton is strictly zero, achieving absolute topological self-locking.

\subsection{Precise Calculation and Emergence of Proton Physical Properties}

\subsubsection{Asymptotic Freedom and Color Confinement: A Fluidic Interpretation}

In the Coupled-Mode Equations (\ref{eq:cme_forward}-\ref{eq:cme_backward}), the three quarks are locked together through their strong interactions with the gluon field $A^\mu_G$:
\begin{itemize}
    \item \textbf{Color Confinement}: Since $\sum \delta_j = 2\pi$, these three complex phases form a closed loop in the complex plane. If one attempts to isolate a single quark ($F_3 \to 0$), the gluon field $A^\mu_G$ can no longer maintain local phase compensation, resulting in massive, uncancelled imaginary dissipation terms. To preserve algebraic self-consistency, the grid exerts a massive restoring force via the logarithmic term $\ln(1+u/u_{th})$ along the spatial gradient (manifesting as a linear confinement potential), preventing wavepacket separation.
    
    \item \textbf{Asymptotic Freedom}: In the central region where the three lobes (quarks) are extremely close to each other (energy density $u_{total} \to \infty$), the derivative of the logarithmic term spontaneously vanishes: $\lim_{u \to \infty} g'(u) \to 0$. The self-focusing force saturates, and the sub-wavepackets behave as free waves, naturally explaining asymptotic freedom.
\end{itemize}

\subsubsection{Precise Calculation of the Proton Charge Radius $R_p$}

Because the proton possesses a three-lobed spatial structure ($C_3$ symmetry), its three-dimensional energy distribution is modulated by the $l=3$ spherical harmonic mode. In classical fluid mechanics, its radial probability distribution function $P(r)$ is expressed as:
\begin{equation}
P(r) = r^2 |R_p(r)|^2 = \mathcal{N}^2 r^8 e^{-2r/\bar{\lambda}_p} \label{eq:proton_radial}
\end{equation}
Finding the boundary of maximum energy concentration (the charge radius) of this composite soliton:
\begin{equation}
\frac{dP(r)}{dr} = 0 \implies 8r^7 e^{-2r/\bar{\lambda}_p} - \frac{2}{\bar{\lambda}_p} r^8 e^{-2r/\bar{\lambda}_p} = 0 \implies R_p = 4 \bar{\lambda}_p \label{eq:radius_deriv}
\end{equation}
The reduced Compton wavelength of the proton is $\bar{\lambda}_p = \frac{\hbar}{m_p c} \approx 0.210309 \text{ fm}$. Substituting this into the equation yields:
\begin{equation}
R_p = 4 \times 0.210309 \text{ fm} \approx \mathbf{0.84124 \text{ fm}} \label{eq:radius_val}
\end{equation}

This theoretical value has a relative error of \textbf{only 0.076\%} compared to the accepted experimental proton charge radius of \textbf{$0.8406 \text{ fm}$}. Of course, in a more rigorous theoretical framework, this specific exponential solution should be viewed as a successful phenomenological morphological alignment under the $l=3$ spherical harmonic symmetry constraint; the true self-consistent wave function profile must be obtained by solving the nonlinear field equation (\ref{eq:field_eq}) numerically for the multi-body system.

\subsection{Proton Anomalous Magnetic Moment and the Emergence of the Self-Locking Phase Scale Law}

\subsubsection{The Universal ``Self-Locking Phase Magnetic Moment Scaling Law''}

We first propose a fundamental physical law within the fluid-mechanical framework of the computational grid—the \textbf{Self-Locking Phase Magnetic Moment Scaling Law}: if a self-trapped electromagnetic soliton has its external deformed boundary radius locked by nonlinear self-focusing to $N$ times its reduced Compton wavelength (i.e., $R = N \bar{\lambda}$), and its boundary flow velocity is constrained by the rigid vacuum light-speed limit ($v_{edge} = c$), then the baseline magnetic moment $\mu_{base}$ generated by its overall circulation strictly satisfies:
\begin{equation}
\mu_{base} = N \cdot \mu_{\text{magneton}} \label{eq:scaling_law}
\end{equation}

\textbf{Analytical Proof}: \\
For a fluid wavepacket carrying charge $e$ and moving with angular velocity $\omega$ on a circular path of radius $R$, the classical formula for its magnetic dipole moment is:
\begin{equation}
\mu = \frac{1}{2} e \omega R^2 \label{eq:classical_mu}
\end{equation}
Since the boundary rotational angular velocity satisfies the rigid vacuum light-speed self-locking condition $\omega = c/R = \frac{c}{N\bar{\lambda}}$, and the characteristic radius is $R = N\bar{\lambda}$, substituting these into the equation yields:
\begin{equation}
\mu_{base} = \frac{1}{2} e \left( \frac{c}{N\bar{\lambda}} \right) (N\bar{\lambda})^2 = N \cdot \left( \frac{1}{2} e c \bar{\lambda} \right) \label{eq:proof_step}
\end{equation}
According to the definition of the nuclear magneton $\mu_{N} \equiv \frac{e\hbar}{2m_p} = \frac{1}{2} e c \bar{\lambda}_p$, this simplifies directly to: $\mu_{base} = N \cdot \mu_N$. \textbf{Q.E.D.}

This scaling law uncovers the \textbf{geometric nature of particle magnetic moments}: it depends entirely on the separation factor $N$ of the soliton's boundary relative to its Compton scale:
\begin{itemize}
    \item \textbf{For the electron} (where the self-focusing radius is $R_e \approx 1\bar{\lambda}_e$): $N=1$, and its baseline magnetic moment is strictly $1 \mu_B$.
    \item \textbf{For the proton} (where the charge radius is derived as $R_p = 4\bar{\lambda}_p$ in Section 2.5): $N=4$, and its baseline circulation magnetic moment is strictly \textbf{$4.000 \mu_N$}.
\end{itemize}

\subsubsection{Binding Magnetic Moment Correction Under Strong Interactions}

In real space, the three quarks are bound together via strong nonlinear self-focusing within the same spatial resonant cavity. As pointed out in the D'Alembertian analysis in Section 2.4, to eliminate the non-aligned residual of fractional phases, the gluon field must maintain an adaptive phase gradient $\nabla_\perp \delta_G(\theta)$.

Since the \textbf{spatial gradient of a phase represents the velocity field of the fluid} ($\mathbf{v} \propto \nabla_\perp \delta_G$), this implies that a \textbf{transverse circulation} maintained by the gluon field must exist between the three quark wavepackets. Because its direction is opposite to the main spin, this circulation naturally manifests as a \textbf{diamagnetic shielding response}. In the GPD framework, this corresponds to the \textbf{binding magnetic moment correction ($\Delta \mu_{binding}$)}.

Algebraically, this binding correction is locked by the diagonal projection coefficient of the underlying grid metric:
\begin{equation}
\Delta \mu_{binding} = \frac{\sqrt{2}+1}{2} \mu_N \approx 1.207 \mu_N \label{eq:binding_corr}
\end{equation}
This algebraic form $\frac{\sqrt{2}+1}{2}$ corresponds to the spatial geometric mean factor generated when the diagonal flow on the grid is projected onto the main spin axis. We propose that at a deeper level, this stems from the \textbf{topological imprint of the cubic lattice on the emergent continuous fluid.}

\subsubsection{Synthesis and Precise Alignment of the Proton Magnetic Moment}

In the fluid topological model of ``Computational Realism,'' since the proton is a rotating composite structure with three lobes, its baseline circulation already fully represents the macroscopic spatial distribution of charge. Therefore, we must discard the traditional quantum mechanical practice of ``superimposing an intrinsic self-energy magnetic moment of a point particle on top of the macroscopic loop current.''

The physical magnetic moment $\mu_p$ presented by the proton to macroscopic observers is obtained by direct algebraic superposition of the baseline circulation magnetic moment and the binding magnetic moment correction:
\begin{equation}
\mu_p = \mu_{base} - \Delta \mu_{binding} = 4.000 \mu_N - 1.207 \mu_N = \mathbf{2.793 \mu_N} \label{eq:mu_p_final}
\end{equation}

\textbf{Comparison with Experiment}:
\begin{itemize}
    \item \textbf{Theoretical value derived from our classical fluid topology}: $\mu_p = \mathbf{2.793 \mu_N}$
    \item \textbf{Accepted experimental value of the proton magnetic moment}: $\mu_{exp} \approx \mathbf{2.792847 \mu_N}$
    \item \textbf{Relative Error}: $\text{Relative Error} < \mathbf{0.01\%}$
\end{itemize}

The anomalous magnetic moment of the proton is not an isolated parameter requiring phenomenological fitting, but rather the result of pure classical topological fluid evolution determined jointly by the \textbf{self-locking scaling law ($4\mu_N$)} and the \textbf{gluon transverse diamagnetic self-shielding ($1.207\mu_N$)}.

\subsection{Analytical Continuation of Microscopic Proton Sources to Complex Topological Manifolds}

In the macroscopic dynamical trial solutions of Sections 2.2 to 2.4, the proton was described as the superposition of three spatially separated fractionalized wavepackets and a gluon phase compensation field. However, from the fluid topological perspective of ``Computational Realism,'' these three phase-singularity sources (modeled as $\delta$-functions) must be singular points of a single, continuous complex potential that is analytic on the complex spatiotemporal manifold. This section performs a rigorous analytical continuation from the discrete lock-phases of quarks to prove that their macroscopic morphology (including spin self-locking and the $3\theta'$ envelope) is a necessary geometric emergence of multi-pole nonlinear interference on the grid.

\subsubsection{Two-Dimensional Topological Vortex Sources of Discrete Fractional Charges}

We project the transverse rotational polarization plane of the proton onto the complex plane $w = x + iy = \rho e^{i\theta}$. The three valence quarks are modeled as wavepackets distributed symmetrically at $120^\circ$ intervals on a characteristic radius $\rho_0$:
\begin{equation}
w_1 = \rho_0, \quad w_2 = \rho_0 e^{i 2\pi/3}, \quad w_3 = \rho_0 e^{i 4\pi/3} \label{eq:complex_roots}
\end{equation}
According to the charge-phase locking rule, they carry fractional topological charges of $q_1 = +2/3, q_2 = +2/3, q_3 = -1/3$, respectively. In incompressible complex fluid dynamics, these three point-like sources are essentially \textbf{fractional two-dimensional topological vortices} located at $w_j$, with a topological charge density satisfying $\rho_q(w) = \sum q_j \delta^{(2)}(w - w_j)$.

\subsubsection{Analytical Continuation of the Complex Potential of Localized Wavepackets}

According to the Cauchy-Riemann constraints, the complex potential $\Omega(w)$ of the two-dimensional irrotational flow field excited by these three discrete topological charges is given by the linear superposition of logarithmic Green's functions:
\begin{equation}
\Omega(w) = \sum_{j=1}^3 q_j \ln(w - w_j) \label{eq:complex_potential_omega}
\end{equation}
Each logarithmic term $\ln(w - w_j)$ introduces a branch point in the complex plane. When the fluid completes a full rotation around $w_j$, its phase undergoes a spontaneous jump of $2\pi q_j$. This irreducible localized topological multi-valuedness provides a purely geometric and absolute barrier for \textbf{``quark color confinement''}.

To obtain a continuous representation of the transverse vector potential, we map the complex potential exponentially, defining the topological generating function $\Psi_{uud}(w) \equiv e^{\Omega(w)}$:
\begin{equation}
\Psi_{uud}(w) = \exp\left( \sum_{j=1}^3 q_j \ln(w - w_j) \right) = \prod_{j=1}^3 (w - w_j)^{q_j} \label{eq:gen_func_prod}
\end{equation}
Substituting the characteristic quark charges yields the strict mathematical expression for the complex topological branch cuts:
\begin{equation}
\mathbf{\Psi_{uud}(w) = (w - w_1)^{2/3} (w - w_2)^{2/3} (w - w_3)^{-1/3}} \label{eq:confinement_branch}
\end{equation}

\subsubsection{Center of Charge Translation and Extraction of the Second-Order Quadrupole Moment}

To extract the macroscopic observable field of the proton (the far-field asymptotic behavior $|w| \gg \rho_0$), we must perform a multipole expansion of Eq. (\ref{eq:confinement_branch}).

\textbf{First-order charge dipole moment and coordinate translation}: \\
The $n$-th order charge multipole moment of the system is defined as $S_n = \sum_{j=1}^3 q_j w_j^n$. Although the geometric coordinates are strictly symmetric ($\sum w_j = 0$), the first-order charge dipole moment is non-zero due to the asymmetric charge distribution:
\begin{equation}
S_1 = \sum_{j=1}^3 q_j w_j = \rho_0 \left[ \frac{2}{3} + \frac{2}{3}e^{i2\pi/3} - \frac{1}{3}e^{i4\pi/3} \right] = \rho_0 e^{i\pi/3} \neq 0 \label{eq:s1_val}
\end{equation}
Since $S_1 \neq 0$, the ``center of charge'' of the proton is shifted from its geometric center. To eliminate the macroscopic first-order dipole phase distortion, we translate the origin of the complex plane to the \textbf{Center of Charge}:
\begin{equation}
w' = w - \rho_0 e^{i\pi/3} \label{eq:coord_translation}
\end{equation}

\textbf{Computation of the second-order charge quadrupole moment}: \\
In the translated center-of-charge coordinate system $w' = \rho' e^{i\theta'}$, the first-order multipole moment is strictly eliminated ($S_1' \equiv 0$). The dominant perturbation of the system is then determined by the second-order quadrupole moment $S_2' = \sum q_j (w_j')^2$. Upon algebraic calculation:
\begin{equation}
S_2' = \sum_{j=1}^3 q_j w_j^2 - 2w_{cc} S_1 + w_{cc}^2 \sum q_j = 2\rho_0^2 e^{-i\pi/3} \neq 0 \label{eq:quadrupole_val}
\end{equation}
Expanding the analytically continued expression (\ref{eq:confinement_branch}) in the far-field limit ($|w'| \gg \rho_0$) of the new coordinate system up to the second-order terms:
\begin{equation}
\Psi_{uud}(w') \approx w' \exp\left( -\frac{S_2'}{2 (w')^2} \right) \approx \rho' e^{i\theta'} - \frac{\rho_0^2 e^{-i\pi/3}}{\rho'} e^{-i\theta'} \label{eq:multipole_approx}
\end{equation}

\subsubsection{Confinement Projection and Fluidic Isomorphism of the Three-Lobed ($3\theta'$) Macroscopic Modulation}

Let us examine the physical meaning of the expansion (\ref{eq:multipole_approx}): it consists of the coherent superposition of the $n=1$ fundamental carrier mode ($\rho' e^{i\theta'}$) and the $n=-1$ quadrupole perturbation sideband mode ($e^{-i\theta'}$).

In the dynamical model of Section 2.4, we had to introduce the \textbf{confinement projection operator $\hat{\mathcal{P}}_{conf}$} to extract the macroscopic observable state. On the complex analytical manifold, this physical operation gains an elegant geometric dual—\textbf{Modulus Normalization} $\frac{\Psi}{|\Psi|}$.  
By extracting the pure phase field, we analyze the phase interference caused by the small perturbation $\delta = \rho_0^2 / \rho'^2$:
\begin{equation}
\frac{\Psi_{uud}(w')}{\left| \Psi_{uud}(w') \right|} \approx \frac{\rho' e^{i\theta'} \left( 1 - \delta e^{-i(2\theta' + \pi/3)} \right)}{\rho' \left( 1 - \delta \cos(2\theta' + \pi/3) \right)} \approx e^{i\theta'} \exp\left[ i \delta \sin\left(2\theta' + \frac{\pi}{3}\right) \right] \label{eq:modulus_norm}
\end{equation}

\textbf{Spontaneous emergence of the $3\theta'$ structure}: \\
When this nonlinear phase field couples to the real-space wave envelope, the high-frequency sideband phase difference excites a composite angular harmonic:
\begin{equation}
e^{i\theta'} \cdot i \sin\left(2\theta' + \frac{\pi}{3}\right) \propto e^{i(3\theta' + \pi/3)} - e^{-i(\theta' + \pi/3)} \label{eq:harmonic_emergence}
\end{equation}
Due to the self-focusing response of the underlying nonlinear field equation (phase-amplitude coupling), this $n=3$ pure phase perturbation spontaneously translates into an amplitude modulation of the real envelope, giving rise to the symmetric three-lobed envelope structure:
\begin{equation}
\text{Macro Envelope} \propto 1 + \varepsilon \cos(3\theta' + \phi) f(\rho') \label{eq:envelope_final_emerge}
\end{equation}

\textbf{Section Summary}: \\
Through rigorous complex analytical continuation and center-of-charge translation, we have demonstrated that \textbf{the second-order quadrupole term of the microscopic complex potential, under modulus normalization (equivalent to confinement projection), spontaneously and uniquely transforms into a three-lobed ($3\theta'$) modulated envelope.} This not only derives the macroscopic proton trial solution of Section 2.3 from pure topological geometry, but also establishes a seamless logical closure between the ``discrete quark components'' and the ``collective topological proton'' within the mathematical framework of the underlying grid.

\section{The Meson Trial Solution and Its Properties}

Unlike baryons, which are absolutely stable and possess independent topologically conserved charges, mesons in the ``Computational Realism'' framework are transient states composed of one quark and one antiquark. This section quantitatively constructs the fluid trial solution of mesons by introducing a unified quark-gluon field and provides a complete resolution to the ``meson scale puzzle.''

\subsection{The Meson Scale Puzzle and the ``Split Ancestry'' Mechanism}

In standard nuclear physics, the physical rest mass of the pion is $m_\pi \approx 139.57 \text{ MeV}$. If we directly applied the single-particle self-locking condition, its characteristic Compton scale would be $\bar{\lambda}_\pi = \frac{\hbar}{m_\pi c} \approx 1.41 \text{ fm}$. If this scale were used as the radius of its radial wave function, the calculated charge radius would far exceed the accepted experimental value of \textbf{$0.659 \text{ fm}$}. This is the famous meson scale puzzle.

This theory proposes that the root of this scale puzzle lies in the following: \textbf{mesons are fundamentally localized wavepackets split and stripped from a baryon ``mother body'' during high-energy collisions or vacuum fluctuations.} \\
Therefore, although the split meson undergoes adaptive phase-locking adjustments (acting as a pseudo-Nambu-Goldstone boson, whose binding energy cancels most of its mass), \textbf{its baseline physical scale at the moment of splitting unconditionally inherits the resolution footprint of its baryon ancestor}:
\begin{equation}
\bar{\lambda}_{footprint} \equiv \bar{\lambda}_p = \frac{\hbar}{m_p c} \approx \mathbf{0.210309 \text{ fm}} \label{eq:meson_scale}
\end{equation}


\subsection{Mathematical Construction of the Meson Trial Solution and Gluon Phase Balancing}

Mesons are composed of one valence quark and one valence antiquark (such as $\pi^+ = u\bar{d}$). In the three-dimensional spacetime fluid, this corresponds to a two-lobed dipole structure satisfying $C_2$ rotational symmetry. We represent the vector potential of the meson as the coherent superposition of quark and gluon components:
\begin{equation}
A_{\pi+}(\rho, \theta, z, t) = \left[ A_{\pi+, Q}(\rho, \theta, z) + A_{\pi+, G}(\rho, \theta, z) \right] e^{i\omega_\pi t} \label{eq:meson_ansatz}
\end{equation}
For the $\pi^+ = u\bar{d}$ system:
\begin{itemize}
    \item The locking phase of the $u$ quark is $\delta_u = 5\pi/6$.
    \item According to the antimatter temporal conjugate rule, the phase of the anti-down quark ($\bar{d}$) is $\delta_{\bar{d}} = -\delta_d + \pi = -\pi/3 + \pi = 2\pi/3$.
\end{itemize}
The sum of their microscopic phases is $\delta_u + \delta_{\bar{d}} = 5\pi/6 + 2\pi/3 = 1.5\pi$.

The macroscopic intrinsic phase of the $\pi^+$ meson must be strictly aligned with its topological eigenvalue of $1.5\pi$ at the boundary. Consequently, the target value for the total interference of the backward waves is projected onto the complex plane as a complex vector pointing in the direction of $1.5\pi$, namely $2 e^{i1.5\pi} = -2i$.

Since $1.5\pi \neq 2\pi$ (not closed), the \textbf{gluon channel $A^\mu_{\pi, G}$ inside the meson must intervene to provide adaptive phase compensation}, balancing the resonant cavity longitudinally and angularly:
\begin{equation}
A_{\pi+, Q} = \sum_{j=1}^2 \left( F_j e^{ik_\pi z} + F_j e^{i\delta_j} e^{-ik_\pi z} \right) \Phi_j(\rho, \theta) \label{eq:meson_quark_comp}
\end{equation}
\begin{equation}
A_{\pi+, G} = F_{G,\pi} e^{ik_\pi z} + F_{G,\pi} e^{i\delta_{G,\pi}(\theta)} e^{-ik_\pi z} \label{eq:meson_gluon_comp}
\end{equation}
Setting the total forward wave amplitude normalized to $+2$, the projection coefficient $C_{G, \pi} \equiv F_{G,\pi} e^{i\delta_{G,\pi}}$ of the gluon field onto the $n=1$ macroscopic fundamental mode satisfies:
\begin{equation}
C_{Q, back} + C_{G, \pi} = -2i \implies C_{G, \pi} = -2i - \left( e^{i5\pi/6} + e^{i2\pi/3} \right) \label{eq:meson_matching_eq}
\end{equation}
\begin{equation}
C_{G, \pi} = -2i - (-1.366 + 1.366i) = 1.366 - 3.366i \approx \mathbf{3.63 e^{-i 1.19 \text{ rad}}} \label{eq:meson_gluon_val}
\end{equation}

\textbf{Physical Interpretation}: \\
This calculation indicates that the gluon field inside the meson carries a phase lag of approximately $-67.9^\circ$ ($-1.19 \text{ rad}$). This adaptive phase compensation successfully absorbs the phase deficit of the two-body system, allowing it to macroscopically reorganize into a lossless $l=2$ spatial standing wave. Furthermore, the relatively large gluon field amplitude ($|C_{G,\pi}| \approx 3.63$) indicates that due to the fractional, asymmetric phase locking, the underlying grid must excite a sufficiently strong localized nonlinear gluon self-focusing field to sustain the self-locking stability of this complex meson soliton.


\subsection{Radial Probability Density Peak Calculation}

Under the spherical harmonic modulation, this $l=2$ state introduces an $r^2$ prefactor to the radial envelope $R_\pi(r)$. Combined with the spatial scale $\bar{\lambda}_p$ inherited from the parent proton, we construct the radial probability distribution function $P_\pi(r)$ of the $\pi$ meson as:
\begin{equation}
P_\pi(r) = r^2 |R_\pi(r)|^2 = \mathcal{N}^2 r^6 e^{-2r/\bar{\lambda}_p} \label{eq:meson_radial}
\end{equation}
Taking the derivative and finding the extremum point (the boundary of maximum energy concentration of this composite soliton):
\begin{equation}
\frac{dP_\pi(r)}{dr} = \mathcal{N}^2 \left( 6r^5 e^{-2r/\bar{\lambda}_p} - \frac{2}{\bar{\lambda}_p} r^6 e^{-2r/\bar{\lambda}_p} \right) = 0 \implies R_{\pi, base} = 3 \bar{\lambda}_p \label{eq:meson_deriv}
\end{equation}
Substituting $\bar{\lambda}_p \approx 0.210309 \text{ fm}$, we calculate the baseline charge radius of the $\pi$ meson before anomalous corrections:
\begin{equation}
R_{\pi, base} = 3 \times 0.210309 \text{ fm} = \mathbf{0.63093 \text{ fm}} \label{eq:meson_radius_val}
\end{equation}

\section{Neutron Instability and the Nucleon Binding Mechanism}

In the preceding sections, we successfully constructed the static topological structures of baryons (protons) and mesons. However, the core mystery of the strong interaction is this: why does the neutron, which is unstable in its free state, gain absolute stability when bound with a proton?

\subsection{Spatiotemporal Topological Morphology and Phase Deficit of the Neutron}

According to the fractional charge-phase mapping rule established in Section 2 ($\delta_q = \frac{1+q}{2}\pi$), the locking phase angles of the three sub-standing waves inside the neutron ($udd$) are $\delta_{u} = 5\pi/6$ and $\delta_{d1} = \delta_{d2} = \pi/3$, respectively. Summing these locking phases algebraically:
\begin{equation}
\sum_{j=1}^3 \delta_j = \delta_u + \delta_{d1} + \delta_{d2} = \frac{5\pi}{6} + \frac{\pi}{3} + \frac{\pi}{3} = \frac{9\pi}{6} = \mathbf{1.5\pi} \label{eq:neutron_phase_sum}
\end{equation}
Since $1.5\pi \neq 2\pi$, the neutron cannot form a closed topological loop in the complex plane ($e^{i1.5\pi} = -i \neq 1$), resulting in continuous grid phase leakage when propagating in vacuum.

\subsubsection{Algebraic Constraints of the Decay Channel}

To eliminate this continuous topological tension and transition into an absolutely stable, closed proton state ($2\pi$), the unstable neutron must emit a wavepacket carrying a specific phase. According to the topological phase conservation law, the phase of the emitted wavepacket $\delta_{emit}$ must satisfy:
\begin{equation}
\delta_n = \delta_p + \delta_{emit} \implies 1.5\pi = 2\pi + \delta_{emit} \implies \delta_{emit} = \mathbf{-0.5\pi = -\frac{\pi}{2}} \label{eq:phase_emit}
\end{equation}

\subsubsection{Feynman-Stueckelberg Alignment}

In complex phase space, emitting a negative phase $-0.5\pi$ into the future time direction is \textbf{strictly equivalent in topology and causality to propagating a positive phase $+1.5\pi$ into the past time direction} (since $e^{-i\pi/2} = e^{i3\pi/2} = -i$). According to our calculation of meson phases in Section 3, the sum of the constituent quark phases of the positive pion ($\pi^+$) is precisely:
\begin{equation}
\delta_{\pi^+} = \delta_u + \delta_{\bar{d}} = \frac{5\pi}{6} + \frac{2\pi}{3} = \mathbf{1.5\pi} \label{eq:pion_phase_match}
\end{equation}
This establishes a precise topological correspondence: \textbf{the emission of a negative pion ($\pi^-$) carrying a $-0.5\pi$ phase is fully isomorphic in time-reversal symmetry to the propagation of a positive pion ($\pi^+$) carrying a $1.5\pi$ intrinsic phase.} This resolves the sign contradiction caused by simple charge addition in neutron decay dynamics.

\subsection{The ``Emission-Capture'' Phase Mechanism of Nucleon Binding}

When a proton and a neutron approach each other to bind into a nucleus (taking the simplest deuteron $d = p + n$ as an example), their wavepackets overlap, carving a \textbf{shared spatial resonant cavity} in the underlying grid.

Within this shared cavity, the unstable neutron utilizes spontaneous spatiotemporal fluctuations to release its excess phase tension by emitting a $-\pi/2$ phase wavepacket (i.e., a negative pion $\pi^-$), allowing its own core to upgrade to the stable proton structure ($2\pi$):
\begin{equation}
1.5\pi - \left( -\frac{\pi}{2} \right) = \mathbf{2\pi} \quad (\text{Neutron } \to \text{ Proton}) \label{eq:n_to_p}
\end{equation}
Due to the physical boundary of the shared resonant cavity, this emitted negative phase wavepacket ($\pi^-$) does not escape to infinity, but is instantaneously \textbf{captured (absorbed)} by the adjacent proton (which carries a full $2\pi$ locking phase). Upon absorbing this $-\pi/2$ phase, the proton's local phase decreases, converting it into a neutron:
\begin{equation}
2\pi + \left( -\frac{\pi}{2} \right) = \mathbf{1.5\pi} \quad (\text{Proton } \to \text{ Neutron}) \label{eq:p_to_n}
\end{equation}

This process establishes a high-frequency, lossless \textbf{``emission-capture'' dynamical loop}:
\begin{equation}
\begin{aligned}
\text{Neutron}(1.5\pi) &\xrightarrow{\text{emit } \pi^-(-\pi/2)} \text{Proton}(2\pi) \\
\text{Proton}(2\pi) &\xrightarrow{\text{capture } \pi^-(-\pi/2)} \text{Neutron}(1.5\pi)
\end{aligned} \label{eq:exchange_loop}
\end{equation}
Through this high-frequency phase exchange, the two nucleons continuously swap identities. On macroscopic timescales, the net phase leakage of the system is strictly zero. By dynamically sharing its phase deficit with its proton partner, the neutron gains \textbf{permanent dynamical stability} within the nucleus.

\subsection{Topological Fluid Proof of Nucleon Pairing Selection Rules}

Using this ``emission-capture'' dynamical model, we can self-consistently address the nucleon pairing selection puzzles in the strong interaction:
\begin{enumerate}
    \item \textbf{Non-binding of the Dineutron ($nn$) system}: If two neutrons attempt to bind, both possess a strong spatiotemporal phase tension of $1.5\pi$, and both tend to emit a $-\pi/2$ phase to upgrade to a proton. However, because \textbf{no $2\pi$ proton state exists within the system to act as a ``capturer''}, the emitted negative phase wavepacket cannot be absorbed and dissipates into the external vacuum, causing the instant nonlinear disintegration of the solitons.
    
    \item \textbf{Non-equilibrium Repulsive Pressure of the Diproton ($pp$) system (Helium-2)}: If two protons attempt to bind, both are already in the absolutely stable, closed $2\pi$ state, and \textbf{no member can act as a ``phase emission source''}. Furthermore, since both protons carry the identical material helicity sequence (manifested as the same $b_1 \to b_0$ delay), their overlapping spatial momentum flows generate an enormous parallel saturation pressure (the classical equivalent of Coulomb repulsion), causing the diproton system to instantly undergo proton emission decay.
    
    \item \textbf{Stable Alignment of the Proton-Neutron (Deuteron, $pn$) system}: One member acts as the emission source (the neutron) and the other as the capturer (the proton), forming a self-consistent dynamical loop. Since the neutron contains two negatively charged $d$ quarks and the proton contains two positively charged $u$ quarks, their overlapping fractional phase wavepackets form localized, counter-rotating gear-like meshings. This complementary flow relieves the local grid compression tension, allowing the attractive nonlinear self-focusing force (the nuclear force) to overcome the residual electrostatic repulsion, establishing the proton-neutron pair as the topological cornerstone of stable nuclei.
\end{enumerate}

\section{Fluid Geometric Origin of Parity Violation and Neutrino Topological Forbiddenness}

In electromagnetic and strong interactions, the dynamics are characterized by coherent standing wave interference inside self-locked topological solitons. However, the weak interaction corresponds to the \textbf{topological tearing and reweaving of the soliton structure}. During such transition states, the real mass term that maintains self-focusing undergoes transient melting ($m \to 0$), decoupling the left-handed and right-handed Weyl spinors. This section utilizes fractional phase-locking analysis in the Weyl representation to derive the fluid-topological origin of maximal parity violation and CP symmetry.

\subsection{Coupling Characteristics of Helicity and Charge Conjugation in the Weyl Representation}

In the complex potential representation of the bilayer grid, the charge $q$ is locked to the phase difference $\delta_q = \frac{1+q}{2}\pi$ between the forward and backward waves. In the massless Weyl representation, the two-component Weyl spinor is expressed as:
\begin{equation}
\Psi = \begin{pmatrix} \psi_L \\ \psi_R \end{pmatrix} \doteq \begin{pmatrix} F \\ F e^{i\delta_q} \end{pmatrix} = \psi_L \begin{pmatrix} 1 \\ e^{i\delta_q} \end{pmatrix} \label{eq:weyl_spinor}
\end{equation}
where $F$ is the spatial real envelope function, $\psi_L$ serves as the reference fluid (with phase set to $0$), and the right-handed component $\psi_R$ carries the phase offset $\delta_q$ locked by the charge.

\subsubsection{Linear Decoupling in the Lepton Sector}

For leptons carrying integer charges (such as the electron, $q = -1$), the locking phase is $\delta_{elec} = 0$. Applying the charge conjugation operator $\hat{C} = -i\sigma_y \mathcal{K}$ (where $\mathcal{K}$ is the complex conjugation operator) in the Weyl space:
\begin{equation}
\Psi_e = \begin{pmatrix} F \\ F \end{pmatrix} \xrightarrow{\hat{C}} \Psi_e^c = \begin{pmatrix} 0 & -1 \\ 1 & 0 \end{pmatrix} \begin{pmatrix} F \\ F \end{pmatrix} = \begin{pmatrix} -F \\ F \end{pmatrix} \label{eq:lepton_c_trans}
\end{equation}
The ratio of the left-handed and right-handed components changes simply from $+1$ to $-1$, which physically represents a simple sign flip. This indicates that in integer-charged systems, helicity and charge conjugation (C-symmetry) are algebraically decoupled.

\subsubsection{Non-Trivial Asymmetric Coupling in Hadrons}

For fractional-charge hadron components (such as the up quark $u$, $q = +2/3$, with phase $\delta_u = 5\pi/6$), the response of the Weyl spinor under C-transformation is:
\begin{equation}
\Psi_u = \begin{pmatrix} F \\ F e^{i\delta_q} \end{pmatrix} \xrightarrow{\hat{C}} \Psi_u^c = \begin{pmatrix} 0 & -1 \\ 1 & 0 \end{pmatrix} \begin{pmatrix} F^* \\ F^* e^{-i\delta_q} \end{pmatrix} = \begin{pmatrix} -F e^{-i\delta_q} \\ F \end{pmatrix} \label{eq:hadron_c_trans}
\end{equation}
The new left-handed component of the antiparticle state becomes $\psi'_L = -F e^{-i\delta_q}$, and the new right-handed component is $\psi'_R = F$. Examining the ratio of these components:
\begin{equation}
\frac{\psi'_R}{\psi'_L} = \frac{F}{-F e^{-i\delta_q}} = -e^{i\delta_q} = e^{i(\delta_q + \pi)} \label{eq:hadron_ratio_shift}
\end{equation}
This demonstrates that for fractional phase-locked systems, the charge conjugation transformation (reversing the time sequence) algebraically forces a $\pi$ phase shift on the system's helicity ratio. This uncovers a fundamental geometric feature of hadron physics: \textbf{helicity symmetry (spatial geometry) and C-symmetry (temporal sequence) are non-trivially and asymmetrically coupled through the fractional phase-locking angles}.

\subsection{The $\pi/2$ Fractional Phase of the Neutrino and Topological Incommensurability}

The aforementioned helicity-C symmetry correlation dictates the topological characteristics of the weak interaction channels. Since weak decays are accompanied by the generation of neutral leptons (neutrinos), the selection rules of the weak interaction vertex are constrained by the intrinsic phase of the neutrino.

Substituting the neutrality condition $q=0$ of the neutrino into the phase-locking relation yields a non-trivial, purely imaginary fractional locking phase on the bilayer grid:
\begin{equation}
\left. \delta_\nu = \frac{1+q}{2}\pi \right|_{q=0} = \frac{\pi}{2} \label{eq:neutrino_lock_phase}
\end{equation}
In the standard Weyl representation $\Psi = \begin{pmatrix} \psi_L \\ \psi_R \end{pmatrix}$, this locking relation directly constrains the neutrino's intrinsic two-component structure:
\begin{equation}
\Psi_\nu = \begin{pmatrix} \psi_L \\ \psi_R \end{pmatrix} = \begin{pmatrix} F \\ F e^{i\delta_\nu} \end{pmatrix} = \begin{pmatrix} F \\ iF \end{pmatrix} \label{eq:neutrino_weyl}
\end{equation}
In the spatiotemporal fluid-dynamical picture, Eq. (\ref{eq:neutrino_weyl}) indicates that the neutrino is a physical entity containing a positive photon component ($\psi_L = F$, purely real) and a negative photon component ($\psi_R = iF$, purely imaginary).

\subsection{Topological Residue Forbiddenness and the Dynamical Origin of Helicity Selection}

Analyzing the propagation behavior of this neutral wavepacket on the discrete computational grid: the vacuum grid, as the topological substrate of physical entities, requires the far-field boundary condition of the physical field to be aligned to integers on the real plane (phase shifts can only be integer multiples of $0$ or $\pi$, corresponding to $\pm 1$). This ensures the integer quantization and closure of the topological charge.

When the transition field in the melted, massless state undergoes information reconstruction and propagates outward:
\begin{enumerate}
    \item \textbf{Left-handed component ($\psi_L = F$)}: As a purely real baseline, its phase is aligned with the reference axis of the grid, allowing the grid to transmit this coherent far-field radiation.
    
    \item \textbf{Right-handed component ($\psi_R = iF$)}: It carries a purely imaginary phase of $\pi/2$ which is orthogonally suspended on the complex plane. If this purely imaginary component is allowed to propagate freely to the far-field, it would permanently excite an unclosed fractional topological defect in the vacuum, generating a \textbf{Fractional Topological Residue}.
\end{enumerate}

Therefore, to preserve the global integer quantization of the topological charge, the vacuum grid strictly forbids purely imaginary wavepackets from radiating to the far-field.

This static topological constraint is dynamically implemented through the \textbf{spin-helicity locking mechanism proved in Section 6}:
\begin{itemize}
    \item For a wavepacket propagating in the $+z$ direction, the system selects the \textbf{negative helicity ($s = -1$)}, which strictly cancels the spatiotemporal wave equation of the left-handed channel (the real flow $\psi_L$) to $0$, allowing it to smoothly convert into a free far-field propagating wave;
    
    \item Meanwhile, the wave equation of the right-handed channel (the imaginary flow $\psi_R$) fails to cancel and degenerates into an evanescent near-field containing a $2\omega$ pressure.
\end{itemize}

This algebraic split constitutes the universal helicity selection rule for weak interaction vertices. In any weak decay process involving leptons (such as $\beta$-decay $e^- \leftrightarrow \nu_e + W^-$):
\begin{itemize}
    \item Left-handed electrons ($e^-_L$) can smoothly undergo coherent conversion with the left-handed neutrinos ($\nu_{e,L}$) that are permitted to radiate to the far-field.
    
    \item The right-handed neutrino component converted from right-handed electrons ($e^-_R$), being a localized evanescent near-field, cannot establish a far-field radiation path.
\end{itemize}

Consequently, the physical cross-section of right-handed electrons participating in the weak interaction is locked to zero. In the Dirac operator representation, this selection rule corresponds to the left-handed projection operator:
\begin{equation}
\hat{P}_L = \begin{pmatrix} 1 & 0 \\ 0 & 0 \end{pmatrix} \equiv \frac{1 - \gamma^5}{2} \label{eq:left_projection}
\end{equation}
This proves that the $(V-A)$ structure of the weak interaction is a dynamical selection rule emergent from spin-helicity locking, designed to prevent the generation of fractional topological residues when the grid processes the $\pi/2$ fractional phase-locked neutrino wavepacket.

\subsection{Phase Deficits in Weak Transition States and the ``Topological Relief Valve'' Mechanism of Neutrino Emission}

To elucidate the physical necessity of neutrino emission in the weak interaction, we analyze the phase conservation and topological closure of the underlying grid.

In classical fluid topology, any weak transition state (such as neutron decay: $n \to p + e^- + \bar{\nu}_e$) is fundamentally the tearing and reweaving of the soliton structure. During this transient process, although the local real mass term melts ($m \to 0$), the rigid constraints of total charge and total phase-locking conservation on the grid remain active.

We compute the topological phase balance of the initial and final states of this decay channel. According to the charge-phase mapping $\delta_q = \frac{1+q}{2}\pi$ in Eq. (\ref{eq:charge_phase}):
\begin{itemize}
    \item \textbf{Initial State (Neutron, total phase $\delta_n$)}:
    \begin{equation}
    \delta_n = \delta_u + \delta_{d1} + \delta_{d2} = \frac{5\pi}{6} + \frac{\pi}{3} + \frac{\pi}{3} = \mathbf{1.5\pi} \label{eq:n_phase_decay}
    \end{equation}
    
    \item \textbf{Final State (Proton + Electron + Antineutrino, total phase $\delta_{final}$)}:
    \begin{enumerate}
        \item The phase of the proton ($uud$): $\delta_p = \frac{5\pi}{6} + \frac{5\pi}{6} + \frac{\pi}{3} = \mathbf{2\pi}$
        \item The phase of the electron ($q=-1$): $\delta_e = \frac{1-1}{2}\pi = \mathbf{0}$
        \item The phase of the antineutrino ($q=0$, which is complex-conjugated and inverted to negative due to antiparticle temporal rotation): $\delta_{\bar{\nu}} = -\delta_\nu = \mathbf{-0.5\pi}$
    \end{enumerate}
\end{itemize}
Summing the phases of the final state algebraically:
\begin{equation}
\delta_{final} = \delta_p + \delta_e + \delta_{\bar{\nu}} = 2\pi + 0 - 0.5\pi = \mathbf{1.5\pi} \label{eq:final_phase_decay}
\end{equation}

\textbf{Algebraic Alignment}:
\begin{equation}
\delta_n \equiv \delta_{final} \implies 1.5\pi \equiv 1.5\pi \label{eq:decay_closure}
\end{equation}

This rigorous algebraic equality (\ref{eq:decay_closure}) uncovers a profound classical wave concept: \textbf{the emission of the neutrino (or antineutrino) in weak decay is the ``topological phase deficit'' that the system must evacuate to allow the remaining proton to reorganize into a closed $2\pi$ structure.}

In grid wave dynamics, ordinary photons (electromagnetic radiation) can only carry real phases (integer multiples of $0$ or $\pi$) through space. The $-0.5\pi$ phase left behind in the decay is a \textbf{purely imaginary phase}. Ordinary photons cannot carry this topological phase; \textbf{hence, the system must excite a specialized, neutral wavepacket whose baseline phase is locked to a purely imaginary value. This wavepacket is the neutrino.}

The neutrino acts as a ``Topological Relief Valve'' in the weak interaction: through its $\pi/2$ orthogonal locking of positive and negative photons, it carries away the purely imaginary phase deficit, enabling the final product (the proton) to collapse into a stable real closed state.

\subsection{Fluid Geometric Reconstruction of CP Symmetry under Charge Conjugation}

The aforementioned complex-phase-based dynamical shielding provides a geometric foundation for deriving CP symmetry (matter being left-handed, antimatter being right-handed).

Applying the charge conjugation operator $\hat{C} = -i\sigma_y \mathcal{K}$ to the neutrino spinor $\Psi_\nu = \begin{pmatrix} F \\ iF \end{pmatrix}$ (where the imaginary part is conjugated due to time reversal, $(iF)^* = -iF$):
\begin{equation}
\Psi_{\bar{\nu}} = \hat{C} \Psi_\nu = \begin{pmatrix} 0 & -1 \\ 1 & 0 \end{pmatrix} \begin{pmatrix} F \\ iF \end{pmatrix}^* = \begin{pmatrix} 0 & -1 \\ 1 & 0 \end{pmatrix} \begin{pmatrix} F \\ -iF \end{pmatrix} = \begin{pmatrix} iF \\ F \end{pmatrix} \doteq \begin{pmatrix} \psi_{L, \bar{\nu}} \\ \psi_{R, \bar{\nu}} \end{pmatrix} \label{eq:antineutrino_c}
\end{equation}
The algebraic result (\ref{eq:antineutrino_c}) indicates that after the charge conjugation transformation:
\begin{enumerate}
    \item \textbf{Left-handed antineutrino component ($\psi_{L,\bar{\nu}} = iF$)}: Becomes purely imaginary, degenerating into an evanescent near-field.
    
    \item \textbf{Right-handed antineutrino component ($\psi_{R,\bar{\nu}} = F$)}: Becomes purely real. By selecting the \textbf{positive helicity ($s = +1$)}, its wave equation strictly cancels to $0$ in the $+z$ direction, allowing it to propagate freely as a far-field coherent wave.
\end{enumerate}


The effective projection operator for antiparticle weak interaction channels spontaneously becomes the right-handed projection operator:
\begin{equation}
\hat{P}_R = \begin{pmatrix} 0 & 0 \\ 0 & 1 \end{pmatrix} \equiv \frac{1 + \gamma^5}{2} \label{eq:right_projection}
\end{equation}

This demonstrates that the \textbf{CP symmetry}—matter choosing left-handed bias and antimatter choosing right-handed bias—possesses a clear fluid topological explanation: \textbf{CP symmetry is the geometric invariance requirement ensuring that the far-field radiating component of a complex wavepacket remains aligned with the real baseline of the vacuum grid under the combined action of time reversal (C) and spatial parity (P), leaving no residual topological defects in the vacuum.}

\section{The Fluid Topological Origin of the Neutrino}

In the Standard Model, the neutrino is introduced phenomenologically as a neutral, extremely light fermion with only left-handed helicity and three-generation flavor oscillations. This section will demonstrate that within the fluid-mechanical framework of the computational grid, the neutrino requires no artificial ``zero-value truncation'' or abstract quantum mixing. Its unique physical properties originate from a single trial solution governed by classical nonlinear geometry: \textbf{a $\pi/2$ phase-locked coupling state of a positive and a negative photon propagating in the same direction.}

\subsection{Coaxial Locking of Positive and Negative Photons}

In the underlying computational grid, charge is the temporal delay of cross-layer replication. We have defined:
\begin{itemize}
    \item \textbf{Negative Charge (negative photon $\gamma^-$)}: Corresponds to $b_1 \to b_0$ copy latency, with its baseline phase behaving as an imaginary phase: $i e^{i(kz-\omega t)}$;
    
    \item \textbf{Positive Charge (positive photon $\gamma^+$)}: Corresponds to $b_0 \to b_1$ copy latency, manifesting as a real phase: $e^{i(kz-\omega t)}$.
\end{itemize}

Since the neutrino is an electrically neutral ($q=0$) fermion, in the fluid dynamics of the grid, it is described as \textbf{a coaxially locked coupling state of a positive photon and a negative photon in the same direction ($+z$ direction)}.

We write its complete two-component Weyl spinor trial solution:
\begin{equation}
\Psi_\nu(\rho, z, t) = \begin{pmatrix} \psi_L \\ \psi_R \end{pmatrix} = \begin{pmatrix} F(\rho, z-ct) e^{i(kz-\omega t)} \\ i F(\rho, z-ct) e^{i(kz-\omega t)} \end{pmatrix} = F(\rho, z-ct) e^{i(kz-\omega t)} \begin{pmatrix} 1 \\ i \end{pmatrix} \label{eq:neutrino_ansatz_final}
\end{equation}

\textbf{Physical Analysis}: \\
This trial solution consists of two equal-amplitude components representing negative and positive photons. They share the same traveling envelope $F(\rho, z-ct)$ and propagation phase $e^{i(kz-\omega t)}$, but maintain a strict $\pi/2$ (imaginary unit $i$) phase difference locked by the grid's layer latency. Because their positive and negative momentum flows are symmetric, the net macroscopic charge of this composite soliton is strictly zero.

\subsection{Algebraic Self-Cancellation of the Equation of Motion and Helicity Selection}

In three-dimensional space, the Weyl equations for massless fermions must contain the spin Pauli matrices $\vec{\sigma}$. For a unidirectional wave propagating along the $+z$ axis, the spatial gradient operator reduces to $\nabla = ik \hat{z}$, yielding $\vec{\sigma} \cdot \nabla = \sigma_z (ik)$.

Within the fluid grid, this two-component wavepacket must possess a unified \textbf{helicity $s$ (the projection of the spin $\sigma_z$ onto the direction of momentum, $s = \pm 1$)}.

Under the Weyl representation, the spatiotemporal evolution operators of the decoupled channels are:
\begin{itemize}
    \item \textbf{Left-handed channel equation (acting on $\psi_L$)}:
    \begin{equation}
    i(\partial_t - c \sigma_z \partial_z)\psi_L = 0 \label{eq:weyl_l}
    \end{equation}
    
    \item \textbf{Right-handed channel equation (acting on $\psi_R$)}:
    \begin{equation}
    i(\partial_t + c \sigma_z \partial_z)\psi_R = 0 \label{eq:weyl_r}
    \end{equation}
\end{itemize}

Substituting the components $\psi_L = F e^{i(kz-\omega t)}$ and $\psi_R = i F e^{i(kz-\omega t)}$ into these equations (where $\partial_t \to -i\omega$, $\partial_z \to ik$, and $\omega = ck$), under the \textbf{negative helicity ($s = -1 \implies \sigma_z = -1$)} configuration:
\begin{enumerate}
    \item \textbf{Perfect self-cancellation of the left-handed channel}:
    \begin{equation}
    i(\partial_t + c\partial_z)\psi_L = i(-i\omega + c(ik)) \psi_L = (\omega - ck)\psi_L = \mathbf{0} \label{eq:cancel_l}
    \end{equation}
    *(The left-handed equation is algebraically and strictly equal to $0$, allowing lossless propagation in vacuum.)*
    
    \item \textbf{Unbalanced obstruction of the right-handed channel}:
    \begin{equation}
    i(\partial_t - c\partial_z)\psi_R = i(-i\omega - c(ik)) \psi_R = (\omega + ck)\psi_R = \mathbf{2\omega \psi_R} \neq 0 \label{eq:non_cancel_r}
    \end{equation}
    *(Failing to cancel, the imaginary right-handed channel produces a non-aligned residual, degenerating into an evanescent near-field.)*
\end{enumerate}

This algebraic split provides a physical screening mechanism: to allow the real baseline component ($\psi_L$) to propagate freely, while preventing the unclosed imaginary negative-photon component ($\psi_R$) from radiating to the far-field, \textbf{the system is dynamically forced to choose the $s = -1$ (negative helicity) configuration}.

This strictly corresponds to the \textbf{left-handed neutrino} of standard particle physics. It demonstrates from classical wave algebra that \textbf{the neutrino is forced to collapse into a left-handed state because of the presence of the imaginary orthogonal channel (the negative photon component). This is the underlying fluid-topological origin of Parity Violation.}

\subsection{D'Alembertian Operator Cancellation}

Because the neutrino consists of a pair of coaxially co-propagating positive and negative photons, it lacks the static Bragg grating produced by the strong interference of forward and backward waves (yielding no rest mass). Its envelope $F(\rho, z-ct)$ propagates at the speed of light $c$.

Applying the four-dimensional D'Alembertian operator to the envelope results in complete algebraic cancellation:
\begin{equation}
\left( \frac{1}{c^2}\partial_t^2 - \partial_z^2 \right) F(\rho, z-ct) \equiv 0 \label{eq:box_cancel}
\end{equation}
The D'Alembertian operator collapses into a purely two-dimensional transverse Laplacian ($\square \to -\nabla_\perp^2$), and the non-local nonlinear field equation simplifies directly to:
\begin{equation}
\exp\left( -a^2 \nabla_\perp^2 \right) F(\rho) = \left[ 1 + \kappa \ln\left(1 + \frac{\epsilon_0 \omega_\nu^2 |F(\rho)|^2}{u_{th}}\right) \right] F(\rho) \label{eq:neutrino_transverse}
\end{equation}
where $a = \bar{\lambda}_c = c/\omega_\nu$.

\textbf{Physical Self-Focusing}: \\
Eq. (\ref{eq:neutrino_transverse}) is a standard non-local two-dimensional nonlinear soliton equation. The left-hand exponential operator represents transverse diffraction (spreading), while the right-hand logarithmic term provides an attractive central pressure (self-focusing). Their balance stabilizes the neutrino transversely into a \textbf{spatiotemporal filament} with a radius on the order of its Compton wavelength.

\subsection{Classical Fluid Analogue of Flavor Oscillations: Transverse Mode Beating}

Having established that the neutrino's tiny mass originates from a minute nonlinear leakage $\epsilon_\nu$ of its orthogonal phase-locking boundary (which generates an extremely weak backward wave and thus an ultra-weak rest mass grating), we can self-consistently explain the ``three-generation flavor'' and ``neutrino oscillation'' phenomena.

In the fluid soliton framework, neutrinos are not different particles with distinct quantum numbers, but rather \textbf{three discrete spatial eigenmodes} of the same two-dimensional transverse nonlinear equation (\ref{eq:neutrino_transverse}) on the grid:
\begin{enumerate}
    \item \textbf{Electron Neutrino ($\nu_e$) $\iff$ Fundamental Mode ($TEM_{00}$ mode)}: A smooth Gaussian profile with no node. It has the highest energy concentration and is topologically the most stable.
    
    \item \textbf{Muon Neutrino ($\nu_\mu$) $\iff$ First-order Mode ($TEM_{10}$ mode)}: Contains one radial node (a central bright spot surrounded by a concentric ring). Maintaining this shape requires larger transverse inertial energy, manifesting as a slightly larger effective mass.
    
    \item \textbf{Tau Neutrino ($\nu_\tau$) $\iff$ Second-order Mode ($TEM_{20}$ mode)}: Contains two radial nodes, possessing the largest effective mass.
\end{enumerate}

\textbf{Fluid Analogue of Oscillations}: \\
As this ``fluid filament'' propagates through vacuum, subtle metric fluctuations of the grid act as perturbations, exciting a superposition of these three transverse spatial eigenmodes.

Because different modes have slightly different transverse propagation constants, their phases slip relative to one another along the propagation path $+z$. In nonlinear wave physics, this is known as \textbf{transverse mode beating}. A macroscopic observer measuring the energy profile along the path will detect a periodic oscillation of the transverse profile between the fundamental and high-order modes.



\end{document}
